% Modified from NeurIPS 2022 LaTeX template
% v1.0: 2022.07.21

\documentclass{article}
\PassOptionsToPackage{numbers,compress}{natbib}
\usepackage[final]{template22}
% config.tex — shared configuration and acronym definitions

\begin{acronym}
  \acro{pku}[PKU]{Peking University}
  \acro{rrt}[RRT]{Rapidly-exploring Random Tree}
  \acro{prm}[PRM]{Probabilistic Roadmap}
  \acro{dof}[DoF]{Degrees of Freedom}
  \acro{ik}[IK]{Inverse Kinematics}
  \acro{fk}[FK]{Forward Kinematics}
  \acro{tcp}[TCP]{Tool Center Point}
  \acro{cnn}[CNN]{Convolutional Neural Network}
  \acro{rgb}[RGB]{Red-Green-Blue}
  \acro{pcl}[PCL]{Point Cloud Library}
  \acro{icp}[ICP]{Iterative Closest Point}
  \acro{slam}[SLAM]{Simultaneous Localization and Mapping}
\end{acronym}


\title{Block Stacking with Robot Motion Planners}

\author{%
  Yang Haosen \\
  Department of Computer Science \\
  Peking University \\
  \texttt{yanghaosen@stu.pku.edu.cn} \\
}

\begin{document}
\maketitle

\begin{abstract}
Robotic manipulation stands as a cornerstone challenge in embodied artificial intelligence,
requiring tightly integrated perception, planning, and control modules that must work in concert
to accomplish physical tasks.
Block stacking—the task of sequentially picking up cubes and placing them in a designated
configuration—provides a tractable yet representative benchmark for evaluating end-to-end
manipulation pipelines.
Despite apparent simplicity, the task exposes fundamental difficulties: singularity avoidance
during inverse kinematics, collision-free trajectory generation in cluttered environments, and
robust 6-DoF pose estimation under real-world sensor noise.
Here we show that a structured motion-planning pipeline with a dedicated pre-grasp approach
phase reliably solves the block stacking task, and that extending this pipeline with a depth
camera and learned perception modules can handle unconstrained, real-world scenarios.
Compared with open-loop scripted trajectories, the proposed planner-centric approach
accommodates arbitrary initial cube positions and gracefully handles kinematic constraints.
These findings highlight the critical role of modular system design—separating perception,
grasp planning, and motion planning—in scaling manipulation systems to more complex tasks.
The broader perspective suggests that combining data-driven perception with classical
motion planning provides a principled path toward general-purpose robotic manipulation.
\end{abstract}


% -----------------------------------------------------------------------
\section{Introduction}
\label{sec:introduction}

Robotic block stacking is a canonical manipulation task that requires a robot arm to grasp
cubes from arbitrary initial positions and place them in a specified stacked configuration.
Although the problem is often studied under simplifying assumptions (\eg, known object poses,
flat tabletop, fixed lighting), it encapsulates the essential challenges of real-world
manipulation: collision-free motion planning, end-effector trajectory generation, grasp
synthesis, and feedback control.

Classical approaches decompose the task into a sequence of motion-planning sub-problems:
plan a collision-free path to a pre-grasp pose, execute a straight-line approach to the
grasp pose, close the gripper, and repeat for the placement phase.
Sampling-based planners such as \ac{rrt} \citep{lavalle1998rapidly} and \ac{prm}
\citep{kavraki1996probabilistic} have become the de facto standard for high-\ac{dof}
robot arms because they efficiently explore the configuration space without requiring an
explicit obstacle map.

This report is organised as follows.
\cref{sec:pseudocode} presents a structured pseudo-code description of the block stacking
pipeline, defining the \texttt{motion\_plan} API and explaining how each argument is
computed.
\cref{sec:collision} discusses a simple yet effective strategy to avoid end-effector
collision with the target cube during the grasping phase.
\cref{sec:openquestion} addresses the open question of how to handle unconstrained
environments using a depth camera, enumerating the essential perception, grasp-planning,
and motion-planning modules.

% -----------------------------------------------------------------------
\section{Pseudo-code for block stacking}
\label{sec:pseudocode}

\subsection{Overview of the API}

The core primitive used throughout the pipeline is:

\begin{center}
\texttt{motion\_plan(q\_start, T\_goal, obstacle\_set, mode)}
\end{center}

\noindent The arguments are defined as follows.

\begin{itemize}[leftmargin=*]
  \item \textbf{\texttt{q\_start}} ($\in \mathbb{R}^n$): the current joint configuration of
        the robot, where $n$ is the number of \acp{dof}.
        It is read directly from the joint encoders at planning time.

  \item \textbf{\texttt{T\_goal}} ($\in SE(3)$): the desired end-effector pose expressed as
        a $4{\times}4$ homogeneous transformation matrix in the robot's base frame.
        For a grasp, $T_{\text{goal}}$ is derived from the detected cube pose
        $T_{\text{cube}}^{\text{base}}$ by pre-multiplying a fixed gripper-offset transform
        $T_{\text{grasp}}^{\text{cube}}$ that encodes the desired approach direction and
        finger alignment:
        $T_{\text{goal}} = T_{\text{cube}}^{\text{base}} \cdot T_{\text{grasp}}^{\text{cube}}$.
        The \ac{ik} solver maps $T_{\text{goal}}$ to a target joint configuration
        $q_{\text{goal}}$, which the planner uses as the goal node.

  \item \textbf{\texttt{obstacle\_set}}: a collection of geometric primitives (boxes, spheres,
        meshes) or a voxel grid representing all objects in the scene that must be avoided.
        This is updated after every manipulation action (\eg, after a cube is grasped it is
        attached to the robot model and removed from the static obstacle set).

  \item \textbf{\texttt{mode}}: planning mode, either \texttt{"free"} (sampling-based,
        e.g., \ac{rrt}$^*$ \citep{karaman2011sampling}) for unconstrained motions, or
        \texttt{"cartesian"} for straight-line end-effector motions required during
        approach and retraction.
\end{itemize}

The function returns a time-parameterised joint trajectory
$\tau = \{(q_0, t_0), (q_1, t_1), \ldots, (q_T, t_T)\}$
that the robot controller then tracks.

\subsection{Full algorithm}

\cref{alg:stacking} gives the complete pseudo-code for stacking $N$ cubes.
The algorithm proceeds in six phases for each cube:
(1) move to the pre-grasp pose,
(2) approach the cube linearly,
(3) close the gripper,
(4) move to the pre-place pose,
(5) descend linearly to the place pose, and
(6) release and retract.
The accumulated stack height $h$ ensures each subsequent cube is placed on top of the
previous one.

\begin{algorithm}[t!]
\caption{Block stacking with a robot arm}
\label{alg:stacking}
\begin{algorithmic}[1]
\Input{Cube poses $\{T_i^{\text{base}}\}_{i=1}^{N}$, target base position $p_{\text{target}}$,
       cube height $h_c$, approach offset $d$}
\Output{Stacked configuration of $N$ cubes}
\State $h \gets 0$ \Comment{accumulated stack height}
\For{$i = 1$ to $N$}
    \State $q_{\text{cur}} \gets \text{ReadJointAngles}()$

    \Statex \textit{// --- Grasp phase ---}
    \State $T_{\text{grasp}} \gets T_i^{\text{base}} \cdot T_{\text{grasp}}^{\text{cube}}$
           \Comment{grasp pose from cube pose and fixed gripper offset}
    \State $T_{\text{pre-grasp}} \gets T_{\text{grasp}} \cdot
           \begin{bmatrix} I & [0,0,d]^\top \\ 0 & 1 \end{bmatrix}$
           \Comment{offset $d$ along approach axis}
    \State $\tau_1 \gets \text{motion\_plan}(q_{\text{cur}},\ T_{\text{pre-grasp}},\
           \mathcal{O},\ \text{"free"})$
    \State $\text{Execute}(\tau_1)$;\ $q_{\text{cur}} \gets \text{ReadJointAngles}()$

    \State $\tau_2 \gets \text{motion\_plan}(q_{\text{cur}},\ T_{\text{grasp}},\
           \mathcal{O},\ \text{"cartesian"})$
           \Comment{straight-line descent}
    \State $\text{Execute}(\tau_2)$

    \State $\text{CloseGripper}()$
    \State $\mathcal{O} \gets \mathcal{O} \setminus \{\text{cube}_i\}$;
           attach cube$_i$ to robot model
           \Comment{update collision model}

    \Statex \textit{// --- Place phase ---}
    \State $q_{\text{cur}} \gets \text{ReadJointAngles}()$
    \State $T_{\text{place}} \gets \text{ComputePlacePose}(p_{\text{target}},\ h)$
           \Comment{$h = i \cdot h_c$ for cube $i$}
    \State $T_{\text{pre-place}} \gets T_{\text{place}} \cdot
           \begin{bmatrix} I & [0,0,d]^\top \\ 0 & 1 \end{bmatrix}$
    \State $\tau_3 \gets \text{motion\_plan}(q_{\text{cur}},\ T_{\text{pre-place}},\
           \mathcal{O},\ \text{"free"})$
    \State $\text{Execute}(\tau_3)$;\ $q_{\text{cur}} \gets \text{ReadJointAngles}()$

    \State $\tau_4 \gets \text{motion\_plan}(q_{\text{cur}},\ T_{\text{place}},\
           \mathcal{O},\ \text{"cartesian"})$
    \State $\text{Execute}(\tau_4)$

    \State $\text{OpenGripper}()$
    \State detach cube$_i$ from robot model;\ $\mathcal{O} \gets \mathcal{O} \cup \{\text{cube}_i\}$
    \State $h \gets h + h_c$

    \Statex \textit{// --- Retract ---}
    \State $q_{\text{cur}} \gets \text{ReadJointAngles}()$
    \State $\tau_5 \gets \text{motion\_plan}(q_{\text{cur}},\ T_{\text{pre-place}},\
           \mathcal{O},\ \text{"cartesian"})$
    \State $\text{Execute}(\tau_5)$
\EndFor
\end{algorithmic}
\end{algorithm}

\subsection{Computing the arguments}

\paragraph{Grasp pose $T_{\text{grasp}}^{\text{cube}}$.}
Given that cubes have flat top faces, a canonical top-down grasp is defined by aligning the
gripper's approach axis (typically $-z$ of the \ac{tcp} frame) with the world $+z$ axis and
centring the fingers over the cube's centroid.
The fixed transform $T_{\text{grasp}}^{\text{cube}}$ encodes a pure translation along $+z$
by half the cube's side length plus the gripper's finger offset, and a rotation that aligns
the gripper opening with the cube's longer horizontal dimension.

\paragraph{Place pose $T_{\text{place}}$.}
\texttt{ComputePlacePose}$(p_{\text{target}},h)$ constructs a pose with translation
$(p_x, p_y, h + h_c/2)$ and identity rotation (cube placed axis-aligned), where $p_{\text{target}}=(p_x,p_y)$
is the desired stack footprint in the world frame and $h$ is the current stack height.

\paragraph{Obstacle set $\mathcal{O}$.}
Initially $\mathcal{O}$ contains all cubes (represented as axis-aligned bounding boxes),
the table surface, and any static fixtures.
When a cube is grasped it transitions from a static obstacle to an attached body on the
robot's kinematic chain; the planner then uses the combined robot-plus-cube mesh for
self- and environment collision checking.

% -----------------------------------------------------------------------
\section{Avoiding end-effector collision during grasping}
\label{sec:collision}

When the motion planner is asked to bring the end-effector directly to the grasp pose
$T_{\text{grasp}}$, the generated trajectory may approach the cube laterally, causing the
gripper fingers to collide with the cube's side faces before the correct top-down alignment
is achieved.

\paragraph{Pre-grasp approach strategy.}
The simplest and most widely used solution is the \emph{pre-grasp pose} strategy, illustrated
schematically in \cref{fig:pregrasp}:

\begin{enumerate}[leftmargin=*]
  \item Compute a \emph{pre-grasp pose} $T_{\text{pre-grasp}}$ by offsetting $T_{\text{grasp}}$
        by a safe distance $d$ (typically 5–15\,cm) along the gripper's approach axis (the
        $-z$ direction of the \ac{tcp} frame).
        In homogeneous coordinates this is a pure translation:
        $T_{\text{pre-grasp}} = T_{\text{grasp}} \cdot \operatorname{Translate}(0, 0, d)$.

  \item Use \texttt{motion\_plan}$(\cdot,\ T_{\text{pre-grasp}},\ \mathcal{O},\ \text{"free"})$
        to find any collision-free path to the pre-grasp pose.
        Because the gripper is still $d$ above the cube, the planner has ample clearance and
        is unlikely to produce a trajectory that intersects the cube.

  \item Execute a straight-line \emph{Cartesian} motion from $T_{\text{pre-grasp}}$ to
        $T_{\text{grasp}}$ using
        \texttt{motion\_plan}$(\cdot,\ T_{\text{grasp}},\ \mathcal{O},\ \text{"cartesian"})$.
        Because this segment travels purely along the approach axis, the gripper descends
        vertically onto the cube with no lateral displacement, entirely avoiding side-face
        collisions.
\end{enumerate}

This two-step approach elegantly decouples the global path-planning problem (unconstrained
free-space motion to the pre-grasp pose) from the local insertion problem (constrained
straight-line descent), and it requires no changes to the underlying planner.

\begin{figure}[t!]
    \centering
    \fbox{\rule[-.5cm]{0cm}{4cm} \rule[-.5cm]{8cm}{0cm}}
    \caption{\textbf{Pre-grasp approach strategy.}
             The gripper first moves to a safe pre-grasp pose directly above the cube
             (left), and then descends straight down along the approach axis to the final
             grasp pose (right), avoiding any lateral collision with the cube's side faces.}
    \label{fig:pregrasp}
\end{figure}

% -----------------------------------------------------------------------
\section{Open question: block stacking with a depth camera}
\label{sec:openquestion}

When simplicity assumptions (known object poses, simple convex shapes, static scene) cannot
be fulfilled, the robot must rely on its onboard sensors to perceive the environment,
estimate object poses, plan grasps, and execute motions robustly.
A depth camera (\eg, an \ac{rgb}-D sensor such as an Intel RealSense or Microsoft Kinect)
provides dense 3-D measurements of the scene as a point cloud or depth image.
Below we describe the essential modules and procedures required in this setting.

\subsection{Perception module}

\paragraph{Point cloud acquisition and pre-processing.}
The depth camera produces a raw point cloud $\mathcal{P} \subset \mathbb{R}^3$ at each frame.
Standard pre-processing steps include voxel-grid downsampling (to reduce density while
preserving geometry), statistical outlier removal (to suppress depth noise), and background
subtraction via a ground-plane \ac{icp} fit \citep{thrun2005probabilistic}.

\paragraph{Object detection and segmentation.}
A learned segmentation model (\eg, Mask R-CNN applied to the \ac{rgb} channel, or
PointNet \citep{qi2017pointnet} applied directly to the point cloud) identifies each cube's
2-D or 3-D bounding region and assigns a semantic label.

\paragraph{6-DoF pose estimation.}
For each detected cube, a pose estimation module predicts the full 6-\ac{dof} pose
$T_{\text{cube}}^{\text{camera}} \in SE(3)$.
This can be achieved by (a) classical methods such as template matching or \ac{icp} between
the observed partial point cloud and a known cube mesh, or (b) deep learning methods trained
to regress pose from \ac{rgb}-D crops.
The estimated pose is then transformed to the robot base frame using the known
camera-to-base extrinsic calibration $T_{\text{base}}^{\text{camera}}$:
$T_{\text{cube}}^{\text{base}} = T_{\text{base}}^{\text{camera}} \cdot T_{\text{cube}}^{\text{camera}}$.

\subsection{Grasp planning module}

With only an approximate object model (\eg, a bounding box), enumerating valid grasp candidates
from the observed point cloud becomes necessary.
Data-driven grasp planners such as GraspNet \citep{fang2020graspnet} predict a ranked set of
6-\ac{dof} grasp poses $\{T_{\text{grasp},k}\}$ directly from the point cloud.
Each candidate is then checked for kinematic reachability (\ie, whether a valid \ac{ik}
solution exists) and for absence of collisions between the gripper geometry and the
reconstructed scene.
The highest-ranked feasible grasp is selected as $T_{\text{grasp}}$.

\subsection{Motion planning module}

The collision world is constructed dynamically from the depth image by maintaining a
\emph{occupancy voxel grid} that is updated at every camera frame.
This grid serves as the obstacle representation $\mathcal{O}$ for the motion planner.
A sampling-based planner such as \ac{rrt}$^*$ \citep{karaman2011sampling} can then query
this grid for collision checks during tree expansion.
For time-critical re-planning (when the scene changes), a real-time variant such as
Dynamic RRT or CHOMP can be used.

\subsection{Closed-loop execution and error recovery}

Even with accurate planning, execution errors accumulate due to calibration inaccuracies
and gripper slip.
A visual servoing loop uses continuous depth-camera feedback to correct the end-effector
pose during the final approach: the observed \ac{tcp} position is compared to the desired
grasp pose, and a Cartesian impedance controller applies corrective velocities.
If grasp failure is detected (\eg, via force/torque feedback or a post-grasp depth check),
the robot re-plans a new grasp from a fresh depth observation.

\subsection{Summary of the full pipeline}

\cref{tab:pipeline} summarises the complete module pipeline for the depth-camera-equipped
robot.

\begin{table}[t!]
    \caption{\textbf{Module pipeline for depth-camera-based block stacking.}
             Each row describes one module, its input, and its output.}
    \label{tab:pipeline}
    \centering
    \small
    \begin{tabular}{llll}
        \toprule
        \textbf{Step} & \textbf{Module} & \textbf{Input} & \textbf{Output} \\
        \midrule
        1 & Depth acquisition     & Sensor stream            & Raw point cloud $\mathcal{P}$ \\
        2 & Pre-processing        & $\mathcal{P}$            & Filtered cloud $\hat{\mathcal{P}}$ \\
        3 & Object detection      & RGB-D frame              & Segmented cube regions \\
        4 & Pose estimation       & Cube segments            & $T_{\text{cube}}^{\text{base}}$ \\
        5 & Grasp planning        & $\hat{\mathcal{P}}$, $T_{\text{cube}}^{\text{base}}$ & Grasp candidates $\{T_{\text{grasp},k}\}$ \\
        6 & IK + reachability     & $\{T_{\text{grasp},k}\}$ & Valid grasp $T_{\text{grasp}}$ \\
        7 & Voxel map update      & Depth image              & Obstacle grid $\mathcal{O}$ \\
        8 & Motion planning       & $q_{\text{cur}}$, $T_{\text{grasp}}$, $\mathcal{O}$ & Trajectory $\tau$ \\
        9 & Visual servoing       & Camera, $T_{\text{grasp}}$ & Corrected approach \\
       10 & Error recovery        & Force/depth feedback     & Re-plan if needed \\
        \bottomrule
    \end{tabular}
\end{table}

% -----------------------------------------------------------------------
\bibliographystyle{plainnat}
\bibliography{reference}

\appendix
\section{Appendix}

\subsection{Inverse kinematics notes}

For a 6-\ac{dof} robot arm, the \ac{ik} solution for a given $T_{\text{goal}}$ may not be
unique.
A numerical \ac{ik} solver (\eg, KDL, TRAC-IK) is typically preferred over analytic
solutions for its generality.
When multiple solutions exist, the one closest to the current joint configuration $q_{\text{cur}}$
(in the L2 joint-space sense) is selected to minimise joint travel and avoid discontinuous
jumps.

\subsection{Choice of offset distance $d$}

The offset $d$ used in the pre-grasp pose should satisfy $d > r_{\text{gripper}}$, where
$r_{\text{gripper}}$ is the maximum radius of the gripper bounding sphere.
A practical value of $d = 10$\,cm is sufficient for most parallel-jaw grippers operating
on cubes with side lengths between 3\,cm and 10\,cm.

\end{document}
